\section{An oracle-explicit constrained investment economy}\label{sec:constrained47}
We now execute the feasible-witness construction in a coupled two-state economy. Capital affects both future production conditions and current investment capacity. This makes feasible repair part of the economic policy rather than an algebraic example appended to a common-action calculation.

Let $x\in[0,1]^2$, $\beta=15/16$, and $0\le a\le\kappa(x)=1/8+(x_1+x_2)/16$. A common innovation $z$ is continuously uniform on $[-1/32,1/32]$, with opposite exposures in the two transitions:
\begin{align}\label{eq:economy47}
 F_1(x,a,z)&=1/16+x_1/2+x_2/8+x_1(1-x_2)/16+a/2+z,\\
 F_2(x,a,z)&=1/16+x_2/2+x_1/8+x_2(1-x_1)/16+a/4-z.\notag
\end{align}
Define the shortage and imbalance cost
\[
 b(x)=(x_1-x_2)^2/4+2(1/2-x_1-x_2)_+^2.
\]
Stage cost is $\sum_{j=1}^2(x_j-5/8)^2+b(x)+pa^2+4a^4$; terminal cost doubles the squared-target terms and retains $b$. The predeclared cells have $T\in\{2,3\}$ and $p\in\{1,4\}$. These are normalized theoretical economies, not estimated parameters or calibrated welfare magnitudes.

\begin{proposition}[Finite primitives and continuous-law queries]\label{prop:primitive47}
The state domain is invariant under~\eqref{eq:economy47}. In the $\ell^1$ state norm, valid primitive constants are
\begin{equation}\label{eq:moduli47}
 C_x=13/4,\quad L_g=9/2,\quad F_x=11/16,\quad
 F_a=3/4,\quad C_a=p/2+1/4,\quad K^A=1/16.
\end{equation}
The uniform tensor state grid with $N$ subdivisions has covering radius $h=1/N$. At each node, $N+1$ equally spaced feasible actions have covering radius at most $k=1/(8N)$. The repair is $\min\{a_i,\kappa(x)\}$ and has zero extra displacement allowance. If a continuation is $L$-Lipschitz, $M$ uniform midpoint cells enclose its original continuous-law expectation after adding the deterministic remainder $L/(32M)$ to the numerical integration error.
\end{proposition}
\begin{proof}
The primitive formulas give $F_1\in[1/32,29/32]$ and $F_2\in[1/32,27/32]$ by interval bounds. The sums of absolute derivatives in each state column are at most $11/16$, and the action derivative has $\ell^1$ norm $3/4$. Differentiating the costs on the two shortage regions gives~\eqref{eq:moduli47}. Capacity is $1/16$-Lipschitz and at most $1/4$, proving the grid and repair assertions. Within a uniform shock cell, the mean distance to its midpoint is $1/(64M)$. Since the transition has shock modulus two in $\ell^1$, the expectation error is at most $L/(32M)$. Thus quadrature does not change the innovation law to a discrete one.
\end{proof}

The two generators are the feasible cone witness and conventional fitted-value iteration in the tensor bilinear basis. Each works backward against its own fitted future, uses the same feasible action net, and evaluates the same continuous law with the preceding remainder. No optimal value or other method's continuation is supplied as a label. Rounded labels, their valid query errors, slopes, node actions and all policy bounds are stored.

The conventional actor selects a nearest state node and repairs its action at the current capacity. It receives the one-sided account rather than only the older separated certificate. For the bilinear interpolant, the Bellman-residual enclosure is
\[
 -e_t-L_th\le f_t-\mathcal T_tf_{t+1}
       \le D_t^Qk+e_t+L_th.
\]
The interpolation weights are nonnegative and reproduce coordinates. Their weighted mean $\ell^1$ distance to the surrounding vertices is at most $1/N$, which proves the stated extension from the nodal enclosure.

\begin{lemma}[Complete cell-corner verification]\label{lem:corners47}
For the bilinear actor, the selected-policy allowance is $e_t+D_t$, where
\[
 D_t=\max_i\sup_{x\in C_i}
       \{y_{t,i}+L_t\|x-x_i\|_1-f_t(x)\}.
\]
Here $C_i$ is the closed nearest-node cell and both adjacent choices are admitted on a boundary. It suffices to evaluate the at most nine vertices obtained by splitting each such cell along its node coordinates. Outward endpoint evaluation supplies a valid upper bound on $D_t$ without an additional verification mesh.
\end{lemma}
\begin{proof}
On each resulting subrectangle the distance term is affine and the continuation is bilinear. An affine function minus a bilinear function attains its maximum at a vertex: first maximize over one coordinate, then over the other. The label's feasible action, repaired at the queried state, has objective at most $y_{t,i}+e_t+L_t\|x-x_i\|_1$ by feasible transport. Taking the complete maximum proves the allowance.
\end{proof}

\subsection{A cap obtained from economic accuracy}
In this class the previously abstract covering, integration and repair operations are finite routines with explicit costs. Let $S_T=\sum_{t<T}\beta^t$, $a_p=C_x+K^AC_a$ and $b=\beta(F_x+K^AF_a)=705/1024$. Define
\[
 \overline L=\max\{L_g,a_p/(1-b)\},\qquad
 \overline D=C_a+\beta F_a\overline L.
\]
The primitive recursion gives $L_t\le\overline L$ and $D_t^Q\le\overline D$ at every horizon. Suppose that, apart from the quadrature remainder, each nonterminal numerical query has error at most $\xi$ and terminal labels have error at most $\xi_T$. With $M=N$, Theorem~\ref{thm:feasible46} implies
\begin{align}\label{eq:explicitcap47}
 \mathcal G_N&\le C_{T,p}/N+2S_T\xi+2\beta^T\xi_T,\\
 C_{T,p}&=S_T\{\overline D/8+2\overline L+\beta\overline L/16\}
                     +2\beta^TL_g.\notag
\end{align}
Thus a dyadic $N\ge2C_{T,p}/\varepsilon$ and $\xi,\xi_T\le\varepsilon/[4(S_T+\beta^T)]$ suffice for loss at most $\varepsilon$. Acquisition and action quantization add the explicit terms of~\eqref{eq:acquiredgap47}; these must receive part of the same tolerance. Fixed binary64, a fixed label lattice and a finite three-rung study are not asserted to achieve every positive $\varepsilon$.

For this two-state, one-action, one-shock implementation, Bellman queries number $T(N+1)^3$. Each has $N$ midpoint evaluations. Direct cone evaluation examines $(N+1)^2$ cones, giving $O(TN^6)$ primitive arithmetic work. Bilinear future evaluation uses a local patch, giving $O(TN^4)$ work for that implementation. Both store $O(TN^2)$ critic labels and actions. The present batched program additionally holds $O(N^3)$ state--action arrays and $O(BN^2)$ temporary cone entries for fixed batch size $B$. Feasible acquisition and repair add their recorded deployment work, rather than being included at zero cost.

With $p_b$-bit arithmetic and multiplication cost $\mathsf M(p_b)$, the corresponding arithmetic upper bounds acquire a factor $\mathsf M(p_b+T+\log N)$, with coefficient conversion, comparison and serialization charged separately. The precision must also meet the query enclosure requirement. These are algorithmic upper bounds, not lower bounds for all conventional methods. Proposition~\ref{prop:representation47} shows that the identical native envelope has the same leading factors as its ReLU realization. No representation label removes the tensor dependence.

\subsection{Execution and acquired deployment}
The fixed ladder is $N=M=4,8,16$, with three isolated processes per method and economic cell. The 24 services store all 72 rungs, including failed attempts. Targets are $4,2,1,1/2$ in the normalized cost units. They test the construction-to-certificate service; a mathematically valid coarse bound is not automatically an economically adequate tolerance. The full attained bounds and work therefore accompany the target counts.

\input{tables/constrained47}

For each of the twelve distinct saved witness policies we also execute the selector on a $33$-by-$33$ measurement grid, with coordinate acquisition radii $0$, $2^{-12}$ and $2^{-8}$. Exact robust repair and downward action spacing $2^{-12}$ give 72 deployment cases. Capacity is bounded on the entire acquisition box; it is not checked only at its midpoint. The binary64 score allowance is uniform by~\eqref{eq:leaf47}; the grid records exercise implementation and repair activity rather than replacing that proof.

\input{tables/acquisition47}
